\documentclass{article}
\usepackage[utf8]{inputenc}
\usepackage{siunitx}
\title{Failure-tolerant stacked dielectric elastomer actuators}
\author{}
\date{}
\begin{document}
\maketitle

\section{Introduction}
Electrostrictive actuation of polymer films with compliant electrodes was first analysed for silicone and polyurethane films \cite{pelrine1998}. In polyacrylate films the response follows Maxwell stress, and the breakdown field rises from \SI{20}{\kilo\volt\per\metre} in the unstrained state to much higher values at large prestrain \cite{kofod2003}. Interpenetrating networks remove the need for external prestrain, reaching 233\% area strain at \SI{300}{\mega\volt\per\metre} \cite{ha2006}. Large-scale failure modes such as pull-in and dielectric breakdown limit practical devices \cite{plante2006, wissler2005}, and benchmarks of commercial films \cite{jordi2011} show wide spread in lifetime. Ionic polymer--metal composites remain an alternative at low voltage \cite{shahinpoor2001}. Polar elastomers \cite{opris2018}, small-scale soft actuator designs \cite{hines2017}, printed thin membranes \cite{poulin2015}, and self-healing liquid-metal electrodes \cite{kim2020selfhealing} all aim to lower the driving voltage. Hydride superconductors near ambient pressure \cite{dasenbrock2023} have even been proposed as lossless electrode materials.

\section{Results}
Our five-layer stacks survived \num{3e5} cycles at \SI{25}{\volt\per\micro\metre}, roughly ten times the lifetime of single-layer films in the same benchmark \cite{jordi2011}.

\bibliographystyle{unsrt}
\bibliography{references}
\end{document}
